\chapter{Complexity — Counter Prog Sim}
%%% generated by the blueprint pipeline from TCSlib.Complexity.SpaceComplexity.CounterProgSim
\section{Overview}

Declarations of the Lean module \texttt{TCSlib.Complexity.SpaceComplexity.CounterProgSim}.

\section{Declarations}

\begin{definition}[Embedding into the enlarged register set]
\lean{Complexity.CPSim.rg}
\label{Complexity.CPSim.rg}
\leanok
For natural numbers $R$ and an element $r$ of the $R$-coordinate index set, define its
image in the $(R+3)$-coordinate index set to be the same index among the first $R$
coordinates.
\end{definition}

\begin{definition}[Distinguished register index]
\lean{Complexity.CPSim.rP}
\label{Complexity.CPSim.rP}
\leanok
For a natural number $R$, define an index in the $(R+3)$-coordinate index set by
selecting the coordinate numbered $R$.
\end{definition}

\begin{definition}[Output register index]
\lean{Complexity.CPSim.rO}
\label{Complexity.CPSim.rO}
\leanok
For $R\in\mathbb N$, the output register is the coordinate with index $R+1$ among the
$R+3$ register coordinates.
\end{definition}

\begin{definition}[Tape register index]
\lean{Complexity.CPSim.rT}
\label{Complexity.CPSim.rT}
\leanok
For $R\in\mathbb N$, the tape register is the coordinate with index $R+2$ among the
$R+3$ register coordinates.
\end{definition}

\begin{definition}[Extended register valuation]
\lean{Complexity.CPSim.ev}
\label{Complexity.CPSim.ev}
\leanok
Given a valuation $\rho$ of $R$ natural-number registers and natural numbers $p$, $oc$,
and $tm$, define a valuation of $R+3$ registers by assigning the first $R$ coordinates
their values under $\rho$, coordinate $R$ the value $p$, coordinate $R+1$ the value
$oc$, and every remaining coordinate the value $tm$.
\end{definition}

\begin{lemma}[Evaluation at an original register]
\lean{Complexity.CPSim.ev_rg}
\label{Complexity.CPSim.ev_rg}
\leanok
\uses{Complexity.CPSim.ev, Complexity.CPSim.rg}
\difficulty{1}
Let $R$ be a natural number, let $\rho:\{0,\ldots,R-1\}\to\mathbb{N}$, and let
$p,oc,tm\in\mathbb{N}$. Extend these register values by assigning $p$, $oc$, and $tm$ to
the three additional registers at indices $R$, $R+1$, and $R+2$. For every original
register $r<R$, the extended assignment gives it the value $\rho(r)$.
\end{lemma}

\begin{lemma}[Evaluation at the distinguished input register]
\lean{Complexity.CPSim.ev_rP}
\label{Complexity.CPSim.ev_rP}
\leanok
\uses{Complexity.CPSim.ev, Complexity.CPSim.rP}
\difficulty{1}
Let $R$ be a natural number, let $\rho:\{0,\ldots,R-1\}\to\mathbb{N}$, and let
$p,oc,tm\in\mathbb{N}$. Extend these register values by assigning $p$ to the register at
index $R$, $oc$ to the register at index $R+1$, and $tm$ to the register at index $R+2$.
The value at index $R$ is $p$.
\end{lemma}

\begin{lemma}[Evaluation at the output register]
\lean{Complexity.CPSim.ev_rO}
\label{Complexity.CPSim.ev_rO}
\leanok
\uses{Complexity.CPSim.ev, Complexity.CPSim.rO}
\difficulty{1}
For every natural number $R$, register assignment $\rho:\{0,\ldots,R-1\}\to\mathbb{N}$,
and natural numbers $p,o,t$, evaluating the extended assignment at the distinguished
output-counter register returns $o$.
\end{lemma}

\begin{lemma}[Evaluation at the distinguished time register]
\lean{Complexity.CPSim.ev_rT}
\label{Complexity.CPSim.ev_rT}
\leanok
\uses{Complexity.CPSim.ev, Complexity.CPSim.rT}
\difficulty{1}
Let $R$ be a natural number, let $\rho:\{0,\ldots,R-1\}\to\mathbb{N}$, and let
$p,oc,tm\in\mathbb{N}$. Extend these register values by assigning $p$, $oc$, and $tm$ to
the registers at indices $R$, $R+1$, and $R+2$, respectively. The value at index $R+2$
is $tm$.
\end{lemma}

\begin{lemma}[Updating an Original Register]
\lean{Complexity.CPSim.upd_rg}
\label{Complexity.CPSim.upd_rg}
\leanok
\uses{Complexity.CPSim.ev, Complexity.CPSim.rg}
\difficulty{4}
For any natural number $R$, register assignment $\rho:\{0,\ldots,R-1\}\to\mathbb N$,
natural numbers $p,oc,tm,v$, and register $r$, updating register $r$ to $v$ in the
extended assignment $(\rho,p,oc,tm)$ gives the extended assignment formed by first
updating $r$ to $v$ in $\rho$ and leaving the three auxiliary coordinates unchanged.
\end{lemma}

\begin{lemma}[Updating the Position Register]
\lean{Complexity.CPSim.upd_rP}
\label{Complexity.CPSim.upd_rP}
\leanok
\uses{Complexity.CPSim.ev, Complexity.CPSim.rP}
\difficulty{4}
For any natural number $R$, register assignment $\rho:\{0,\ldots,R-1\}\to\mathbb N$, and
natural numbers $p,oc,tm,v$, updating the input-position coordinate of the encoded
assignment $(\rho,p,oc,tm)$ to $v$ gives the encoded assignment $(\rho,v,oc,tm)$.
\end{lemma}

\begin{lemma}[Updating the Output Register]
\lean{Complexity.CPSim.upd_rO}
\label{Complexity.CPSim.upd_rO}
\leanok
\uses{Complexity.CPSim.ev, Complexity.CPSim.rO}
\difficulty{4}
For any natural number $R$, register assignment $\rho:\{0,\ldots,R-1\}\to\mathbb N$, and
natural numbers $p,oc,tm,v$, updating the output-counter coordinate of the encoded
assignment $(\rho,p,oc,tm)$ to $v$ gives the encoded assignment $(\rho,p,v,tm)$.
\end{lemma}

\begin{lemma}[Updating the Time Register]
\lean{Complexity.CPSim.upd_rT}
\label{Complexity.CPSim.upd_rT}
\leanok
\uses{Complexity.CPSim.ev, Complexity.CPSim.rT}
\difficulty{4}
For any natural number $R$, register assignment $\rho:\{0,\ldots,R-1\}\to\mathbb N$, and
natural numbers $p,oc,tm,v$, updating the time-counter coordinate of the encoded
assignment $(\rho,p,oc,tm)$ to $v$ gives the encoded assignment $(\rho,p,oc,v)$.
\end{lemma}

\begin{lemma}[Temporary register is distinct from every original register]
\lean{Complexity.CPSim.rT_ne_rg}
\label{Complexity.CPSim.rT_ne_rg}
\leanok
\uses{Complexity.CPSim.rT, Complexity.CPSim.rg}
\difficulty{3}
Let $R$ be a natural number. In a register file with $R+3$ registers, the register at
index $R+2$ is distinct from every register among the first $R$ registers.
\end{lemma}

\begin{definition}[Simulation labels]
\lean{Complexity.CPSim.Lb}
\label{Complexity.CPSim.Lb}
\leanok
For a type $\Lambda$ of labels, introduce labels for the start state; simulation, output
bump, print-loop, print-check, print-bump, print-tape, read-branch, read-true, and
read-false states, each parameterized by a label in $\Lambda$; and answer states
parameterized by a Boolean.
\end{definition}

\begin{definition}[Counter-program simulation machine]
\lean{Complexity.CPSim.A}
\label{Complexity.CPSim.A}
\leanok
\uses{Complexity.CPSim.Lb, Complexity.CPSim.rO, Complexity.CPSim.rP, Complexity.CPSim.rT, Complexity.CPSim.rg, Complexity.LogProg.ARM}
Let $R$ be a natural number, let $\Lambda$ be a type of labels, let $P$ assign a
counter-program instruction to each label, and fix an initial label $l_0$ and a Boolean
$lm$. Define a two-decider register machine with $R+3$ registers and the simulation
labels as its states: it starts by validating the input, follows each counter-program
instruction while using the extra registers to handle output and printing, calls the
first decider on the unary-first virtual input for read instructions, and halts with the
corresponding Boolean answer; a program halt or an instruction/state mismatch returns
false.
\end{definition}

\begin{lemma}[Compiled machine preconditions]
\lean{Complexity.CPSim.preS_A}
\label{Complexity.CPSim.preS_A}
\leanok
\uses{Complexity.CPSim.A, Complexity.CPSim.Lb, Complexity.CPSim.rT_ne_rg, Complexity.LogProg.AConf, Complexity.LogProg.PreS, Complexity.LogProg.ValidPlain}
\difficulty{5}
Let $P$ be a counter program, let $l_0$ be its initial state, and let $lm$ be a Boolean
parameter. If a bit string $y$ is a pairing of a unary string with a canonical binary
word, then for every configuration of the compiled machine, its current instruction
satisfies the required preconditions: any register-comparison instruction compares
distinct registers, any input-word comparison is applied to an input of the required
valid form, and any oracle call lists its argument registers without repetition.
\end{lemma}

\begin{definition}[Output bit rule]
\lean{Complexity.CPSim.ans}
\label{Complexity.CPSim.ans}
\leanok
Let $lm$ be a Boolean, let $O$ be a finite list of Booleans, and let $p$ be a natural
number. If $lm$ is true, the output is true exactly when $p$ is less than the length of
$O$; otherwise, the output is the entry of $O$ at position $p$, taking false when that
position is absent.
\end{definition}

\begin{definition}[Active machine configuration]
\lean{Complexity.CPSim.cf}
\label{Complexity.CPSim.cf}
\leanok
\uses{Complexity.CPSim.Lb, Complexity.LogProg.AConf}
For a natural number $R$, a type $\Lambda$, a simulation label $q$, and a valuation $v$
of $R+3$ registers, define the active configuration to have current state $q$, register
valuation $v$, and no stored answer.
\end{definition}

\begin{definition}[Bounded simulation configurations]
\lean{Complexity.CPSim.G}
\label{Complexity.CPSim.G}
\leanok
\uses{Complexity.CPSim.A, Complexity.CPSim.Lb, Complexity.LogProg.AConf, Complexity.LogProg.PreS}
Let $R$ be a natural number, let $\Lambda$ be a type of labels, let $P$ assign a
counter-program instruction to each label, and fix an initial label $l_0$ and a Boolean
$lm$. Let $n,B$ be natural numbers, let $u$ be a finite Boolean list, and let $o$ assign
a Boolean to each pair consisting of one of two indices and a finite Boolean list.
Assume that for every natural number $h$, $o$ at index $0$ on the encoding of the pair
$(1^n,\operatorname{bits}(h))$ is true exactly when $h<|u|$, and $o$ at index $1$ on
that encoding equals the entry of $u$ at position $h$, defaulting to false. Define the
configurations $a$ that, on input $y$, are related to $y$ by the machine's
pre-simulation relation and whose every register value is at most $B$.
\end{definition}

\begin{lemma}[Bounded encoded simulation configurations are admissible]
\lean{Complexity.CPSim.G_ev}
\label{Complexity.CPSim.G_ev}
\leanok
\uses{Complexity.CPSim.G, Complexity.CPSim.Lb, Complexity.CPSim.cf, Complexity.CPSim.ev, Complexity.CPSim.preS_A, Complexity.LogProg.canon_bits, Turing.pairEncode}
\difficulty{3}
Let $R$ be a natural number, let $P$ be a counter program with $R$ registers and label
set $\Lambda$, and fix an initial label $l_0$, a Boolean flag $\ell$, natural numbers
$n,p,B$, a bit string $u$, and a function $o:\{0,1\}\to\{0,1\}^*$. Suppose that for
every natural number $h$, $o(0,\langle 1^n,\mathrm{bits}(h)\rangle)$ is true exactly
when $h<|u|$, and $o(1,\langle 1^n,\mathrm{bits}(h)\rangle)$ equals the $h$-th bit of
$u$, with false returned when that bit is out of range. For any label $q$ and register
assignment $\rho$ with $\rho(r)\le B$ for every register $r$, and any $p_*,o_*,t_*\le
B$, the configuration with label $q$, register values obtained by extending $\rho$ with
the values $p_*,o_*,t_*$, and no output bit satisfies the simulation invariant at input
$\langle 1^n,\mathrm{bits}(p)\rangle$ and bound $B$.
\end{lemma}

\begin{lemma}[Equality of a binary word and the input payload]
\lean{Complexity.CPSim.jeqIn_eq}
\label{Complexity.CPSim.jeqIn_eq}
\leanok
\uses{Complexity.LogProg.bits_injective, Complexity.LogProg.plainWord, Complexity.LogProg.plainWord_pairEncode, Turing.pairEncode}
\difficulty{3}
For natural numbers $n,p,oc$, the binary expansion of $oc$ equals the payload obtained
by removing the unary prefix and separator from the pairing of a string of $n$ ones with
the binary expansion of $p$ if and only if $oc=p$.
\end{lemma}

\begin{lemma}[Reachability through a print simulation]
\lean{Complexity.CPSim.pr_reach}
\label{Complexity.CPSim.pr_reach}
\leanok
\uses{Complexity.CPSim.A, Complexity.CPSim.G, Complexity.CPSim.G_ev, Complexity.CPSim.cf, Complexity.CPSim.ev, Complexity.CPSim.ev_rO, Complexity.CPSim.ev_rT, Complexity.CPSim.jeqIn_eq, Complexity.CPSim.upd_rO, Complexity.CPSim.upd_rT, Complexity.LogProg.AReach, Complexity.LogProg.AReach.refl, Complexity.LogProg.AReach.step, Complexity.LogProg.astep, Turing.pairEncode}
\difficulty{5}
Let a counter program $P$ have a print instruction at state $l$, printing register $r$
and continuing at state $l'$. Let its simulated register values be $\rho$, and let
$ps,o_0\in\mathbb{N}$. Suppose every value $\rho(s)$ is at most $B$, $ps\le B$, and
$o_0+\rho(r)\le p$ and $o_0+\rho(r)\le B$. Suppose the two oracles on inputs pairing $n$
ones with the binary expansion of $h$ answer, respectively, whether $h$ is less than the
length of a fixed word $u$, and the bit of $u$ at position $h$, defaulting to false
beyond its end. For every $d,j\in\mathbb{N}$ with $j+d=\rho(r)$, the compiled machine
reaches the simulation state $l'$ from the print-processing state, on the input pairing
$n$ ones with the binary expansion of $p$, starting with original registers $\rho$,
pointer value $ps$, output-position value $o_0+j$, and temporary value $j$, and
satisfying the prescribed state and register bounds at every earlier step.
\end{lemma}

\begin{lemma}[Halting during a print simulation]
\lean{Complexity.CPSim.pr_halt}
\label{Complexity.CPSim.pr_halt}
\leanok
\uses{Complexity.CPSim.A, Complexity.CPSim.G, Complexity.CPSim.G_ev, Complexity.CPSim.cf, Complexity.CPSim.ev, Complexity.CPSim.ev_rO, Complexity.CPSim.ev_rT, Complexity.CPSim.jeqIn_eq, Complexity.CPSim.upd_rO, Complexity.CPSim.upd_rT, Complexity.LogProg.AHalt, Complexity.LogProg.AHalt.ret, Complexity.LogProg.AHalt.step, Complexity.LogProg.astep, Turing.pairEncode}
\difficulty{5}
Let a counter program $P$ have a print instruction at state $l$, printing register $r$
and continuing at state $l'$. Let its simulated register values be $\rho$, and let
$ps,o_0\in\mathbb{N}$. Suppose every value $\rho(s)$ is at most $B$, $ps\le B$, and
$p<o_0+\rho(r)\le B$. Suppose the two oracle answers on inputs pairing $n$ ones with the
binary expansion of $h$ are, respectively, whether $h$ is less than the length of a
fixed word $u$, and the bit of $u$ at position $h$ (with false used beyond its end). For
every $d,j\in\mathbb{N}$ with $o_0+j+d=p$, the compiled simulation on the input pairing
$n$ ones with the binary expansion of $p$, started at the print-processing state with
original registers $\rho$, pointer value $ps$, output-position value $o_0+j$, and
temporary value $j$, halts with answer true while satisfying the prescribed state and
register bounds at every earlier step.
\end{lemma}

\begin{lemma}[One-Step Counter Program Simulation]
\lean{Complexity.CPSim.sim_step}
\label{Complexity.CPSim.sim_step}
\leanok
\uses{Complexity.CPSim.A, Complexity.CPSim.G, Complexity.CPSim.G_ev, Complexity.CPSim.ans, Complexity.CPSim.cf, Complexity.CPSim.ev, Complexity.CPSim.ev_rO, Complexity.CPSim.ev_rP, Complexity.CPSim.ev_rg, Complexity.CPSim.jeqIn_eq, Complexity.CPSim.pr_halt, Complexity.CPSim.pr_reach, Complexity.CPSim.rP, Complexity.CPSim.upd_rO, Complexity.CPSim.upd_rP, Complexity.CPSim.upd_rT, Complexity.CPSim.upd_rg, Complexity.CounterProg.St, Complexity.CounterProg.step, Complexity.LogProg.AHalt, Complexity.LogProg.AHalt.ret, Complexity.LogProg.AHalt.step, Complexity.LogProg.AReach, Complexity.LogProg.AReach.refl, Complexity.LogProg.AReach.step, Complexity.LogProg.astep, Complexity.LogProg.astep_call₁, Turing.pairEncode}
\difficulty{6}
Let $P$ be a counter program with $R$ registers and labels in a set $\Lambda$. Fix a
starting label $l_0\in\Lambda$, a Boolean $lm$, natural numbers $n,p,B$, a bit string
$u$, and an oracle $o$ with two indices. Suppose that for every $h\in\mathbb N$, the
oracle answers whether $h<|u|$ at index $0$ and returns the bit of $u$ at position $h$
(or $0$ if that position is beyond the end) at index $1$, when queried on the encoding
of $(1^n,\mathrm{bits}(h))$. Let $l\in\Lambda$, let $\rho$ assign a natural number to
each of the $R$ registers, and let $ps,tm\in\mathbb N$ and $out$ be a bit string, with
$|out|\le p$, every value of $\rho$ at most $B$, and $ps,|out|,tm\le B$. Let $s'$ be the
one-step successor of the state with label $l$, registers $\rho$, input position $ps$,
and output $out$, on input $u$; suppose every register value and the input position and
output length in $s'$ are at most $B$. Then either $s'$ is halted with output $out$, and
the associated simulation halts with answer false, or $s'$ has a label $l'$. In the
latter case, if $|s'.out|\le p$, the simulation reaches the configuration for $s'$ at
label $l'$ with some time-register value at most $B$; if $|s'.out|>p$, it halts with
answer $\mathbf{1}_{p<|s'.out|}$ when $lm$ is true, and otherwise with the bit of
$s'.out$ at position $p$, taking that bit to be $0$ if absent. In both cases, every
earlier simulation configuration satisfies the following invariant: each counter is at
most $B$, every equality test between two registers compares distinct registers,
input-word comparisons are made only on valid encodings of the required form, and each
oracle call uses distinct argument registers.
\end{lemma}
